\section{Discussion}

The current artifacts support three practical observations. First, post-alignment discrepancy is not uniform: severity and affected\_versions are where deterministic factual-conflict candidates remain most visible, while published is often a temporal or representation issue, references is frequently incomplete, and cwe\_ids is mostly equivalent with a smaller incomplete/conflict tail. Second, the RQ2 diagnostic shows why a binary same/different view is too coarse: most binary-different field instances are assigned to non-conflict categories by the deterministic baseline. Third, the diagnostic evidence-source profile is uneven by source class and field: severity has broad fetched-text availability, while affected\_versions diagnostic samples include a larger number of GitHub and vendor evidence URLs with lower fetch success.

The key implication is operational rather than purely taxonomic. A downstream consumer should not react to every binary mismatch with source selection. Equivalent and representation differences mainly call for normalization, incomplete records call for augmentation or coverage analysis, temporal differences call for timeline interpretation, and only residual factual conflicts call for evidence-constrained adjudication. This routing view is the main method claim that can be defended before gold labels exist: post-alignment fusion needs a discrepancy type before it needs a winner.

The CWE full-impact audit provides the strongest current positive RQ2 candidate signal. After a post-hoc evidence-enhanced pass over the original priority rows, its direction favors treating fully related ancestor/descendant assignments as representation differences on 10 of 14 primary-seed-disjoint strict rows. Additional evidence resolves four of the original six undecided rows, including two OS-command cases, but it also expands the concrete boundary to four factual conflicts despite complete taxonomy connectivity. Two mechanism-specific rows remain unresolved. Evidence therefore improves candidate coverage from 11/17 to 15/17 without making the rule automatic: taxonomy ancestry still cannot establish that the specific CWE applies to the CVE. The post-hoc selection, uneven snapshot availability, same-model reviewers, and absent human signoff make this a prioritized human-audit result rather than a production-switch result.

The references audit exposes a complementary construct boundary. Exact redirects and preserved advisory or bounty identifiers yield 32/32 strict agreements across four rule families, but every encoded-line row splits the reviewers. Both interpretations are coherent: an encoded \texttt{\%23L} suffix can be intended as a line anchor on one file, yet it is also literal path text in the observed request and returns 404. The narrower audited profile avoids choosing the interpretation that best improves the candidate metric. Its exact 32-row impact alignment is useful development evidence, but only a human-defined resource-identity contract can decide whether malformed line references should be normalized, repaired, or retained as distinct failed references.

The three-row residual audit shows the same boundary at case scale. Source code can make a local weakness mechanism concrete, as the unhandled PyTorch Lightning request-body access does for CWE-248, without turning that development judgment into gold. Official taxonomy metadata can also invalidate an abstraction level for vulnerability mapping, as with CWE-840, while still failing to identify the correct concrete weakness from a narrow patch. Finally, the K3s malformed references produce different labels under two coherent resource identities. More evidence is therefore not interchangeable with an approved construct: code can support a mechanism, but taxonomy and URL identity still require an explicit review contract.

The staged frontier adds an operational stopping argument. A third same-model pass resolves 64\% of its selected rows, but two evidence-enriched passes resolve only 16\% of the next selected remainder and 8\% per reviewer-row decision. These selected-stage ratios are not comparable accuracy estimates, yet they show that generic same-family escalation has reached diminishing candidate yield under the current contract. Because the fixed target is still missed and later targeted diagnostics cannot be promoted, the defensible next investment is construct approval plus independent real-person review, or a later frozen snapshot, rather than another vote on the revealed rows.

The same artifacts also show why the paper must stay cautious. RQ2 now has a large non-human typing-stability diagnostic, but still not a gold-label agreement study: its two pass roles share one model/config, comprise many ephemeral sessions, and cannot compare candidate profiles because all predictions are identical. In this non-human diagnostic, the post-unseal analysis shows that the severity result largely compares a label-only baseline with a richer score/vector/version review construct, while affected\_versions exposes 25 concrete unbounded-record projection losses. The resulting 1,121/1,147 post-hoc fit is explicitly in-sample, and the cross-protocol gate then rejects it: severity moves by -23 and -7 rows on the older primary/review slices but +165 on fresh strict labels, while affected inputs are not preserved consistently. This reversal makes protocol drift and input-projection sensitivity a stronger finding than candidate accuracy. Severity remains non-human, the complete 1,250-row human packet is blank, and the corpus has no human-approved field contract yet.

The later isolated snapshot removes the earlier profile-identifiability failure but not the validation boundary. Three CWE rows differ between current and the taxonomy candidate, yet only one receives strict dual-Codex consensus in the sealed main review. Frozen source evidence then makes both fresh reviewers agree on the candidate's granularity-only direction for all three cases, including the two previously non-strict rows. This strengthens the case-level mechanism explanation, not the profile metric: the rows are selected after the differences are known, and the reviewers still share one model family. The two rejected protocol attempts also show why exact output interfaces belong to the evidence contract: malformed confidence or path representations must fail closed rather than be silently normalized into agreement. More importantly, the strict event-time universe is empty: all 250 reviewed rows are snapshot-external rather than event-time-external. The result therefore supplies three prioritized CWE cases and a reproducible availability finding, not a generalization or method-gain estimate. It also illustrates why a later download is insufficient by itself; temporal validation must bind source event timestamps, frozen profiles, and untouched records simultaneously.

The 50-row CWE audit provides a stronger field-control check but reaches the same claim boundary. Hiding the three differences among 47 controls preserves their candidate direction and yields 48/49 candidate agreement versus 45/49 current agreement, while exposing one evidence disagreement and one profiles-equal miss. This reduces the most direct case-selection cue, but it does not undo that CWE was chosen after profile unsealing and that all reviewers use the same model family. Its failed protocol versions are also substantive findings. Literal set inclusion does not by itself imply semantic incompleteness when the extra CWE may describe a different mechanism, and a citation is not auditable unless its quoted text is a literal frozen substring before the decision is accepted. The V3 pre-write checks make those assumptions executable. They improve auditability, not external correctness, so the sealed +1/231 comparison remains the only main profile metric.

The final 16-row evidence stage reinforces the stopping decision across fields. Every selected row has a successfully retrieved non-empty source record, yet only both references and two affected\_versions rows become strict; severity remains 0/2 and affected\_versions 2/12. The bottleneck is therefore not HTTP availability alone. It is whether the supplied records establish the intended CVSS comparison, URL resource identity, or complete product/package/range relation under a shared contract. Reaching 238/250 staged coverage while failing the 0.40 selected-row resolution gate is not a near-pass to be rescued by threshold relaxation. It is evidence that another correlated vote over the same revealed cohort is unlikely to resolve the remaining construct choices.

The label-free envelope sharpens the practical interpretation further. Because 98.8\% of the sealed predictions are identical, no human labeling of this cohort can produce more than a 1.2-point absolute accuracy difference between current and the CWE profile. The six profiles reduce to two prediction vectors, and every cross-vector pair differs on the same three rows. Even the most favorable possible gold assignment gives an exact two-sided McNemar p-value of 0.25, so no labeling of this cohort can support rejection at alpha 0.05. Six same-direction correctness-discordant rows are only the theoretical minimum for any rejection; exact 0.80 power requires 49, 20, or 12 effective discordant rows under candidate-win probabilities 0.70, 0.80, or 0.90. This label-free diagnostic is not evidence that either profile is accurate, and those power values are conditional assumptions rather than a sample-size commitment. The three differential rows determine the paired sign but cannot validate the taxonomy over the other 247 rows or establish absolute performance and construct error. A future confirmatory cohort should therefore require event-time eligibility, enough naturally occurring sealed profile differences, and an a priori power calculation tied to a human-approved outcome contract.

The full eligible-universe census changes the planning interpretation of that calculation. The six profiles are all distinct over the full revealed universe, but only 34 of 29,740 field instances differ: 29 CWE rows and five references rows under the broader references profile, reduced to three under the audited profile. Thus the sample's 3/250 rate should not be extrapolated as stationary; it comes from a current-status-stratified allocation, and the full prediction changes are concentrated in rare rule-trigger strata. A future strict event-time evaluation should first census sealed predictions over its eligible universe. It can then freeze a disagreement-enriched review layer, including all differences when feasible, for paired method comparison, while retaining a separately probability-sampled control layer for absolute accuracy and prevalence. CVE-level grouping and at most one primary comparison row per CVE should prevent pseudo-replication. Gold labels and reviewer identities must remain hidden from prediction census and selection. On the current revealed snapshot, references profiles have only five or three possible current-versus-candidate discordant CVEs, below the six-row extremal significance minimum, whereas CWE and combined profiles have 29--34 prediction differences but unknown correctness discordance. This is design information, not power or method evidence.

The second label-free acquisition makes the temporal bottleneck observable rather than treating the empty strict tier as a generic waiting condition. NVD adds 39 records published after the profile freeze, while the reviewed GHSA snapshot adds none; therefore no bilateral post-freeze match exists and the frozen field views remain unchanged. This supports a source-aware readiness gate: collection time alone cannot start confirmatory review, and unilateral source growth should not be counted toward a bilateral event-time denominator. The correct action is to retain the sealed profiles, rerun the same label-free availability audit after GHSA updates, and freeze a cohort only when the predeclared minimum is met. This monitoring result is not a temporal validation result because it contains no eligible row and no correctness label.

Auditing the complete revealed difference set sharpens, but does not remove, that boundary. The 25-to-1 CWE direction shows that the taxonomy profile's changed rows are usually compatible with the same-model evidence contract, while the one current-direction row and three unresolved rows show that ancestry is not sufficient as an automatic rule. Conversely, the references result shows why field status alone is coarser than resource identity: reviewers can agree that a target GHSA alias is shared yet disagree on whether NVD and third-party pages belong to that same underlying document group. Requiring complete partition identity exposes this disagreement rather than granting a candidate win from one selected pair. Future evaluation should therefore freeze the resource ontology and adjudication unit with real reviewers before profile comparison, not derive them from whichever partition improves the metric.

The third blind pass does not remove this limitation. Its fixed majority rule resolves 66 of the 103 non-human A/B non-strict rows, but misses both advancement thresholds and leaves 37 non-human candidate rows, including 28 affected\_versions cases. Within these non-human outputs, triple uncertainty on 17 rows and three-way qualified-label splits on seven show two distinct bottlenecks: some inputs lack enough evidence for a determinate decision, while others expose taxonomy or field-contract ambiguity even when the model is confident. A fourth same-family vote would add correlated evidence rather than resolve either construct. The next useful step is evidence enhancement and real-person semantic adjudication, not threshold relaxation.

The evidence-secondary result distinguishes generic evidence availability from construct-complete evidence. Frozen advisory, commit, issue, and release-page text is available for 28 of the 37 non-human rows and supports six strict D/E candidates outside affected\_versions. It resolves no non-human affected\_versions row: both reviewers still abstain together on many cases because text fragments do not establish complete product/package identity, all range endpoints, branch/backport scope, or a total component-to-product release mapping. The shared 19 non-human uncertain decisions make exact agreement a misleading success measure. For this field, the next evidence unit must be a project-specific release, manifest, lockfile, or vendor mapping graph under a human-approved range contract; adding more generic URLs or correlated votes does not address the missing edge.

The disjoint non-human development calibrations make the construct boundary more specific. They support a candidate distinction between same-CVSS-version base-metric conflict and cross-version representation, and they repeat the missing-score and one-sided-unbounded rules on disjoint rows. They also show that high aggregate agreement is insufficient: v2 has perfect exact agreement, yet its only non-strict row is a shared abstention. Official component evidence alone does not resolve that abstention, and the current-coordinate-only diagnostic abstains under its frozen contract. The historical non-human graph shows what additional structure is needed: product-build dependencies map product versions to legacy artifact versions, source identity binds the migration edge, and only then can prerelease sets be ordered. That graph yields a post-hoc non-human \texttt{incomplete} candidate for the XWiki case. A generic fail-closed schema then binds all eight structurally selected equal-boundary cases across three ecosystems, showing that package identity, component membership, and an explicit artifact alias can be represented by one evidence contract. The non-equal audit reveals that released-version set equality is not the same construct as syntactic interval containment. Freezing the assistant's extensional choice and applying it to two InLong components makes that alternative executable, but the representation candidate still conflicts with both reviewers. The full-aligned-data stress test then shows a separate boundary: heterogeneous package claims cannot be unioned merely because registries enumerate both sets. Parallel distributions, dependency constraints, and coordinated components require different evidence edges, and none of the initial NuGet, PyPI, or crates.io rows has a total product-release map. The post-no-go Deno result narrows this claim: a dependency range is insufficient, but an exact committed build lock can recover the product edge and expose five runtime-only affected releases. This is evidence for an edge-specific mapping rule, not for treating all lockfiles as authoritative or for resolving the field's extensional/intensional choice. Real people must still choose the range semantics and approve acceptable lockfile, build-manifest, or vendor mappings for each edge class.

The Mattermost family audit shows that project-specific evidence is necessary but still not sufficient. Exact tag manifests establish the current module identity for every fixed token, and the pseudo-version timestamps identify unique commits, yet most product tags are on incomparable backport branches. Treating timestamp order as release ancestry would manufacture a total mapping that the commit graph does not support. The historical module path adds a second unresolved union edge. The defensible outcome is therefore 0/2 abstention, not a representation or incompleteness label. Together with Deno, this isolates two different edge classes: exact build locks can recover one product-to-component sequence, while repository identity plus a pseudo commit does not recover a multi-branch release projection.

The EVE family makes the failure more specific. A GHSA package string can name a repository that contains several nested Go modules and non-Go build components without being a module itself. In addition, a syntactically valid pseudo version identifies a commit boundary but does not guarantee that the commit modifies the vulnerable component: one EVE pseudo changes a filesystem setup script outside \path{pkg/vtpm}. The public advisory API and repository advisory UI also expose different levels of version detail. A content-addressed commit graph can remove API-rate and reproducibility barriers, but it cannot turn those interface, component-identity, or branch-divergence gaps into a total product set. The 0/2 result therefore supports field-specific abstention and evidence-interface preservation, not a claim that the API, UI, NVD, or GHSA is factually wrong.

Hutool isolates a different condition: a registry can expose a coordinated aggregate/component release line without supplying a human-approved meaning for an unbounded affected range. Equal Maven catalogs make the release tokens comparable, and the anchor POM/JAR evidence prevents a name-only product-to-component projection. Under the declared snapshot-extensional semantics, both NVD sets are strict subsets of the GHSA component unions. That mechanism result is useful, but it does not establish that the current 209-release registry snapshot is the intended historical universe or that every aggregate JAR was inspected. Since the mapping shape was known before v1, the defensible consequence is a new-cohort requirement, not promotion of the two revealed \texttt{incomplete} candidates.

Applying the frozen mechanism to six CVE-exposure-disjoint Hutool rows separates mechanism reuse from repeated judgment on the original two cases. The six projections include both strict-subset and equal-set outcomes, so the contract does not mechanically assign every eligible row to \texttt{incomplete}. This strengthens the evidence that the cached token correspondence is executable across more than the development pair. It still does not establish temporal generalization or accuracy: cohort availability was discovered in the same aligned snapshot, labels remain Codex development candidates, and promotion is disabled.

The affected\_versions field remains a special case because version-system mismatches can still produce diagnostic conflict signals. OSV/GHSA-style affected records separate package identity and range events, and recent affected-version work shows that this field remains difficult even when tools use multiple forms of advisory evidence \cite{OSVSchema2026,GHSARepo2026,Chen2025AffectedVersions}. The current baseline has been tightened and NVD \texttt{vulnerable=false} CPE matches have been filtered, but the residual 651 factual-conflict candidates still need manual spot checks before they can be treated as stable field labels. The repair changed only 10 full-corpus classifications and no frozen-cohort decision, but that stability is an observed input-impact result, not evidence that package identity or range semantics are solved.

The development evidence-aware pass for affected\_versions supports an evidence-availability statement, a silver-label distribution, and a silver-only baseline comparison. The v1 CVE-disjoint holdout adds a stronger leakage-control check: its cohort, evidence cache, blind reviewer input, and predictions were separately frozen, and the predictions predate both decision files. This improves auditability, but it does not convert Codex decisions into human gold. Exact joint consensus covers only 35/100 rows, and the reviewer agreement values indicate substantial unresolved ambiguity.

More importantly, the v1 post-unseal diagnostic changes how the affected\_versions result should be interpreted. The preregistered all-strict diagnostic score rewards a method both for recognizing non-conflict cases and for selecting a source within true factual conflicts. In this diagnostic, token methods score mainly through the former because representation discrepancies tend to map to \texttt{both}; on the 16 strict factual conflicts, no proposed method exceeds fixed GHSA preference or recency. V1 therefore exposes the need to separate discrepancy typing from conditional source adjudication.

V2 preregisters that separation on a new cohort, so its negative result is more informative. In this non-human diagnostic, the type candidate is correct on its three determinate predictions but abstains on 38 of 41 strict type rows; counting abstention as incorrect makes its full accuracy substantially lower than both full-coverage comparators. On the nine non-human strict-source rows, the branch source head reaches 2/9, while fixed NVD preference reaches 6/9. Thus endpoint separation fixes the interpretation problem but does not rescue method performance. The defensible finding is a protocol correction plus a failed, highly selective method candidate, not an affected\_versions adjudicator.

The post-v2 diagnostics sharpen that conclusion rather than overturning it. A more permissive task-separated type candidate and an authority-filtered source candidate do not improve consistently across Phase D, v1, and v2. A leave-one-cohort-out classifier improves the pooled type count but regresses below the legacy comparator on held-out Phase D, while its source endpoint is below the branch baseline both pooled and on Phase D. This cohort sensitivity is why we do not spend another sealed cohort on the current candidate family. The next methodological requirement is a human-audited source-authority contract and version semantics that first demonstrate cross-cohort stability on development data.

The evidence audit also separates reviewer independence from evidence independence. Five of nine v2 strict-source rows use the same single cited URL for both reviewers, and only two have primary/ecosystem evidence cited by both. Agreement under a shared evidence snapshot can therefore reflect a common evidence bottleneck. Future source labels must record source ownership and distinguish upstream maintainer, ecosystem database, disclosure, database record, and secondary evidence before a strict source decision is accepted.

The sensitivity diagnostic reinforces that caution. Severity predictions change for fewer samples across evidence-score thresholds, while affected\_versions predictions change for most samples when the minimum version-token support threshold is raised. That instability is not a failure by itself, but it shows that token presence is a brittle proxy for source support unless package identity and version-range semantics are modeled explicitly.

The case sketches reinforce the same boundary. The severity examples show both the value and the weakness of a simple evidence-score rule: exact score/vector matches can recover source support, but they can also miss lower-salience evidence for the other side. The affected\_versions examples show why this field needs extra care around package identity, patch-level granularity, and ecosystem-specific version schemes before a conflict detector can be treated as reliable.

The silver error-mode diagnostic gives a concrete audit target for the next iteration. Severity mismatches are often source-split differences where the silver label marks both-source support but the baseline marks only one source, or the reverse. Affected\_versions mismatches are more structurally diverse, and the current 100-sample diagnostic highlights many no-overlap package-name cases. This does not validate the silver labels as gold, but it does identify where a future human audit and semantic version-range adjudicator should focus first.

This positioning also constrains the related-work claim. The contribution should not be framed as superseding vulnerability text discrepancy mining, general truth discovery, or affected-version analysis. Its defensible space is narrower: a post-CVE-alignment workflow that makes structured-field disagreement auditable and routes only the residual conflict cases to source-support analysis with abstention.

The main discussion point should therefore be selective and modest: the current evidence-aware results are diagnostics within non-human label setups. The strongest supported contribution is the auditable type-first pipeline, including explicit abstention, sealed prediction artifacts, the identified task-mixing failure mode, evidence-dependence audit, and a follow-up protocol that evaluates the two tasks separately. Any broader claim about affected\_versions method effectiveness or general vulnerability-field fusion must wait for human gold. A further untouched cohort is warranted only after a revised candidate first passes the documented cross-cohort development gate.
